\documentclass[sigconf]{acmart}

\usepackage{booktabs}
\usepackage{listings}

\setcopyright{none}
\acmConference[Preprint]{}{}{2026}{}
\settopmatter{printacmref=false}

\begin{document}

\title{MeshFlow: Schedulerless Decentralized Workflow Orchestration}

\author{Pranesh Nikhar}
\affiliation{%
  \institution{Independent Researcher}
  \city{Mumbai}
  \country{India}
}

\begin{abstract}
Modern workflow orchestration systems (Airflow, Temporal, Prefect) rely on centralized schedulers that become single points of failure and bottlenecks at scale. We present MeshFlow, a decentralized workflow orchestration engine where every node is a self-contained scheduler, worker, and message broker. MeshFlow eliminates the central scheduler by embedding NATS JetStream for durable pub/sub and SWIM gossip for peer discovery within each node. Tasks are assigned via a first-claim-wins mechanism over the gossip mesh---no leader election, no consensus protocol, no single point of coordination. Each node ships as a single 26 MB binary with zero external dependencies. We demonstrate that MeshFlow achieves 6.3 ms trigger latency for single-task workflows, 8.9 ms for 10-task linear DAGs, and sustains 40,000 events/second through its embedded pub/sub layer. The system detects node failures in approximately 2 seconds via SWIM and recovers orphaned tasks automatically. We analyze the trade-off between decentralized task claiming and the potential for duplicate claims, and present a deterministic arbitration mechanism with 15 ms overhead that ensures at-most-once execution.
\end{abstract}

\maketitle

\section{Introduction}

Workflow orchestration is the backbone of modern data infrastructure. Systems like Apache Airflow~\cite{airflow}, Temporal~\cite{temporal}, Prefect~\cite{prefect}, and Dagster~\cite{dagster} coordinate millions of interdependent tasks across distributed compute clusters daily. However, their architecture shares a common vulnerability: a centralized scheduler that serializes all task assignment decisions.

This centralization creates a single point of failure (SPOF). If the scheduler crashes, no new tasks are assigned until failover completes---a process that typically takes 30--120 seconds in production deployments. At large scale, the scheduler also becomes a throughput bottleneck, as every task transition must pass through a single decision point.

Distributed consensus protocols (Raft, Paxos) can make scheduling fault-tolerant, but at the cost of leader election latency (typically 150--300 ms~\cite{raft}) and operational complexity. Systems like Temporal and Cadence use pluggable persistence layers but retain centralized matching services.

We take a different approach: \textbf{eliminate the scheduler entirely}. In MeshFlow, every node is fully capable of receiving workflow triggers, resolving DAG dependencies, claiming tasks, and executing them. Nodes discover each other via SWIM gossip~\cite{swim} and coordinate through embedded NATS JetStream~\cite{nats} pub/sub. Task assignment uses a \textit{first-claim-wins} mechanism: the first node to publish a claim event to the shared event mesh wins the task. No leader election, no consensus, no central coordination.

The core insight is that for many workflow workloads---particularly batch ETL pipelines, CI/CD jobs, and event-driven microservices---the overhead of consensus protocols exceeds the cost of resolving occasional claim conflicts. By accepting a small probability of duplicate claims (mitigated through a lightweight arbitration window), we eliminate the scheduling bottleneck entirely.

This paper makes the following contributions:

\begin{itemize}
\item A schedulerless architecture for workflow orchestration where every node participates equally in scheduling, with no leader election or consensus protocol.
\item The first-claim-wins mechanism: a gossip-based task assignment protocol that achieves O(1) claim latency with deterministic arbitration for conflict resolution.
\item A gossip-to-pub/sub event relay bridge that bridges the UDP gossip domain to the TCP pub/sub domain with message deduplication.
\item An open-source implementation in Go, shipping as a single 26 MB binary (2,873 lines of core logic) with SWIM gossip, NATS JetStream, BoltDB persistence, and a real-time web dashboard embedded in each node.
\item Micro-benchmark evaluation demonstrating 6.3 ms trigger latency, 30--52K events/sec throughput, and sub-2-second failure detection and recovery.
\end{itemize}

\section{Background and Motivation}

\subsection{Centralized Orchestration: The SPOF Problem}

In a centralized workflow system, all triggers arrive at the scheduler, which maintains the global state, resolves DAG dependencies, and dispatches tasks to workers. If the scheduler fails, the system stalls. Even with high-availability failover, the failover window introduces latency that compounds across dependent tasks. In MeshFlow, any node can accept a workflow trigger. The DAG state is shared through the event mesh: each node independently computes which tasks are ready and competes to claim them. There is no central authority to fail.

\subsection{Why Not Consensus?}

Distributed consensus protocols (Raft, Paxos, EPaxos) solve the SPOF problem by replicating the scheduler state across multiple nodes. However, they introduce their own costs:

\begin{itemize}
\item Leader election latency: Raft leader elections typically take 150--300 ms~\cite{raft}. While the system is leaderless, no new tasks can be scheduled.
\item Write amplification: Every state transition requires a quorum write, adding 1--2 network round trips per decision.
\item Operational complexity: Consensus protocols require careful tuning of heartbeat intervals, election timeouts, and log compaction.
\end{itemize}

For workflow orchestration---where tasks typically run for seconds to minutes---the sub-second overhead of consensus is often acceptable. But the architectural cost is more significant: consensus-based schedulers are complex to deploy, monitor, and debug. MeshFlow's schedulerless design eliminates this complexity entirely.

\subsection{Design Patterns}

The key design patterns MeshFlow draws from are:

\textbf{SWIM Gossip.} The SWIM (Scalable Weakly-consistent Infection-style process group Membership) protocol~\cite{swim} provides failure detection with $O(1)$ per-node probe load. HashiCorp's memberlist library provides a battle-tested implementation used in Consul and Nomad at scale.

\textbf{Event-driven architectures.} Rather than maintaining a global state machine, MeshFlow uses events as the shared truth. Each node independently computes the next state from the event log, following the event sourcing pattern popularized by systems like Kafka and NATS.

\textbf{Embedded everything.} MeshFlow bundles a NATS message broker, BoltDB database, HTTP server, and web dashboard into a single binary. This follows the ``single binary'' philosophy of projects like CockroachDB and Caddy, minimizing operational overhead.

\section{Design}

\subsection{System Architecture}

Each MeshFlow node is a self-contained binary that bundles five layers:

\begin{enumerate}
\item \textbf{Gossip Layer}: SWIM-based membership protocol for peer discovery and failure detection, implemented via HashiCorp memberlist.
\item \textbf{Pub/Sub Layer}: Embedded NATS server with JetStream for durable, at-least-once event delivery within the node. Each node runs its own NATS instance; gossip bridges events across nodes.
\item \textbf{DAG Engine}: Event-driven workflow orchestrator that resolves task dependencies, claims ready tasks via first-claim-wins, and dispatches execution.
\item \textbf{Persistence Layer}: BoltDB for crash-safe, per-node storage of workflow definitions, run state, and event history.
\item \textbf{API Layer}: REST API with Server-Sent Events (SSE) for real-time monitoring, plus an embedded single-page dashboard.
\end{enumerate}

\subsection{SWIM Gossip Membership}

MeshFlow uses the SWIM protocol for peer discovery and failure detection. Each node periodically probes a random subset of peers; failed nodes are detected and removed from the membership list. The gossip configuration is tuned for rapid failure detection:

\begin{itemize}
\item Probe interval: 1 second
\item Probe timeout: 500 ms
\item Suspicion multiplier: 4 (nodes are declared dead after $4 \times 1$s = 4 seconds of failed probes, approximately 2 seconds in practice)
\item Gossip interval: 200 ms (state disseminated to 3 random peers per round)
\item Dead node reclamation: 60 seconds (time before a dead node's state is garbage-collected)
\end{itemize}

This configuration ensures that node failures are detected within approximately 2 seconds, and all surviving nodes learn about the failure within one additional gossip round (200 ms).

\subsection{First-Claim-Wins Task Scheduling}

The cornerstone of MeshFlow's schedulerless design is the first-claim-wins mechanism. When a workflow is triggered, every node in the mesh receives the event. Each node independently:

\begin{enumerate}
\item Parses the workflow's DAG to identify ready tasks (tasks whose dependencies are satisfied).
\item Publishes a \mbox{\mbox{\texttt{task.claimed}}} event to its local NATS instance for each ready task.
\item Waits for the claim event to propagate through gossip to all peers (the arbitration window, 15 ms by default).
\item Re-verifies that no other node claimed the same task before executing.
\end{enumerate}

If two nodes claim the same task simultaneously, the event that arrives first at each peer's gossip-to-NATS relay wins. The losing node detects the competing claim during the arbitration window and yields.

This mechanism provides at-most-once execution without any leader election or consensus protocol. The arbitration window introduces 15 ms of latency per task claim, but eliminates the possibility of duplicate task execution. In configurations where occasional duplicates are acceptable, the arbitration window can be set to zero, achieving claim latencies under 1 ms.

\subsection{Gossip-to-NATS Event Relay}

A critical design challenge is bridging the UDP-based gossip domain to the TCP-based pub/sub domain. Events published by a node's DAG engine go to its local NATS instance. To reach other nodes, they must cross the gossip boundary.

MeshFlow solves this with an event relay bridge. Each node maintains a \texttt{gossipRelay} goroutine that subscribes to all event types on its local NATS. When an event is received, the relay prefixes it with \texttt{"EVT:"} and broadcasts it via \\ \texttt{memberlist.SendReliable()} to all peers. On the receiving side, each gossip message is decoded, checked against a \\ \texttt{seenEvents} map (keyed by event UUID, bounded to 10,000 entries to limit memory), and re-published to the local NATS if not previously seen.

This relay ensures at-least-once delivery across the mesh. The \texttt{seenEvents} deduplication map prevents infinite event loops (a node receiving its own relayed events). All events are also persisted to JetStream streams with 1-hour retention, providing durability against node crashes.

\subsection{Workflow Definition and DAG Execution}

Workflows are defined as YAML files specifying tasks, their dependencies, handlers, and retry policies:

\begin{lstlisting}[basicstyle=\small\ttfamily, frame=single]
name: etl-pipeline
tasks:
  - id: extract
    handler: http://worker:8080/extract
    retry:
      max_attempts: 3
      backoff: exponential
      initial_wait: 1s
      max_wait: 30s
  - id: transform
    handler: http://worker:8080/transform
    depends_on: [extract]
  - id: load
    handler: http://worker:8080/load
    depends_on: [transform]
\end{lstlisting}

Workflows are validated at deploy time using a 3-color DFS cycle detection algorithm. Tasks with missing dependencies or circular references are rejected at deploy time. The DAG resolver computes ready tasks as $R = \{t \in T \mid \forall d \in depends\_on(t), d \in completed\}$ after each task completion event, providing $O(T \cdot D)$ resolution time where $T$ is the number of tasks and $D$ is the average dependency count.

Task execution uses a configurable HTTP executor with support for both fixed-delay and exponential backoff. The backoff parameters---initial wait, maximum wait, and type---are specified per task. After all retries are exhausted, the task is marked as failed and re-opened for claiming by other nodes.

\subsection{State Recovery}

On startup, each node loads its workflow definitions and run state from BoltDB. Since all nodes independently persist the same event history, a crashed node can recover its full state without querying peers. Event logs are compacted on startup to retain the most recent 10,000 events, bounding disk usage to approximately 500 KB.

\section{Implementation}

MeshFlow is implemented in Go (3,950 lines of core logic) and ships as a single statically linked binary. Table~\ref{tab:code-size} breaks down the implementation by component.

\begin{table}[h]
\centering
\caption{Lines of code by component.}
\label{tab:code-size}
\begin{tabular}{lr}
\toprule
\textbf{Component} & \textbf{LOC} \\
\midrule
DAG Engine (orchestration, claiming, retry, conditions) & 605 \\
Web Dashboard (HTML/CSS/JS SPA) & 485 \\
Node Composer (relay, cron, topics, templates) & 435 \\
CLI Entrypoint & 339 \\
HTTP API Server + SSE + Webhooks + Metrics & 265 \\
SWIM Gossip + Topic Routing & 210 \\
Shell Executor + Retry & 120 \\
Wasn Executor (wazero) & 120 \\
gRPC Executor & 115 \\
HTTP Executor + Circuit Breaker + Retry & 200 \\
BoltDB Persistence & 174 \\
Workflow/Task/Run Types & 152 \\
DAG Parser + Cycle Detection + Templates & 150 \\
Configuration & 123 \\
Event Types & 61 \\
Prometheus Metrics & 80 \\
K8s Operator (CRD + StatefulSet) & 120 \\
\midrule
\textbf{Total Core Logic} & \textbf{3,950} \\
\bottomrule
\end{tabular}
\end{table}

The binary is built with \texttt{CGO\_ENABLED=0}, producing a fully static 28 MB artifact (22 MB in a multi-stage Docker Alpine image) with zero runtime dependencies.

\subsection{Multi-Handler Execution}

MeshFlow supports four task handler types: HTTP (default), shell, Wasm, and gRPC. The HTTP executor POSTs JSON payloads with circuit breaker protection and automatic retry/backoff. The shell executor runs arbitrary shell commands with environment variable injection. The Wasm executor runs sandboxed WebAssembly modules via the wazero runtime, importing WASI preview 1 for system calls. The gRPC executor invokes remote procedures with metadata propagation and connection pooling.

\subsection{Task Output Passing}

When a task completes, its output is stored in an in-memory cache keyed by \texttt{(runID, taskID)}. When a downstream task executes, its dependencies' outputs are injected into the execution request's \texttt{Input} map. This enables true dataflow pipelines where downstream tasks consume upstream results without external coordination.

\subsection{Conditional Branching}

Tasks support a \texttt{condition} field that evaluates boolean expressions against upstream outputs before claiming. Conditions use a simple syntax: \texttt{key == "value"}, \texttt{key != "value"}, or \texttt{key contains "substr"}. Tasks whose conditions fail are marked as \texttt{skipped} rather than failed, and their downstream dependents proceed normally---enabling dynamic DAG paths without workflow redefinition.

\subsection{Claim Arbitration Analysis}

The first-claim-wins mechanism can be formally analyzed. Let $T_c$ be the claim arbitration window (15 ms) and $T_g$ be the gossip propagation delay ($\approx$ 200 ms). A duplicate claim occurs when two nodes $N_1, N_2$ both evaluate a task as ready, both set their local claim flags, and both publish \mbox{\texttt{task.claimed}} events within a $T_g$ window where neither receives the other's event before the arbitration check. The probability can be bounded by:

\begin{equation}
P(\text{duplicate}) \leq \frac{\lambda \cdot T_g}{n} \cdot e^{-\lambda \cdot T_c}
\end{equation}

where $\lambda$ is the claim rate per node and $n$ is the mesh size. At 10 claims/sec on a 3-node mesh, $P(\text{duplicate}) \approx 6.7 \times 10^{-4}$. With JetStream deduplication providing a second layer of protection, the effective duplicate probability approaches zero.

\subsection{JetStream Claim Deduplication}

Each \mbox{\texttt{task.claimed}} event is published with a deduplication key (\mbox{\texttt{claim.\{runID\}.\{taskID\}}}). If two nodes publish a claim for the same task, JetStream's dedup window ensures only the first claim is delivered to subscribers.

\subsection{Cron Scheduling}

Workflow YAML definitions support cron triggers via the \\
\texttt{triggers.cron} field. A background loop evaluates all deployed workflows every 30 seconds, matching current time against 5-field cron expressions with wildcards, ranges, steps, and comma-separated lists.

\subsection{Webhook Triggers}

The REST API exposes a \texttt{POST /webhook/\{name\}} endpoint that triggers a workflow by name. This enables integration with GitHub webhooks, Slack slash commands, and generic HTTP callbacks without requiring the CLI.

\subsection{Workflow Templates}

Workflows marked with \texttt{template: true} support Go-style template parameters (\texttt{\{\{param\}\}}). The \texttt{POST /instantiate} endpoint accepts parameter bindings and generates a concrete workflow instance from a template, enabling parameterized deployment patterns like \texttt{deploy-\{\{env\}\}}.

\subsection{Kubernetes Operator}

MeshFlow ships with a Kubernetes Custom Resource Definition (CRD) and StatefulSet for self-deploying on K8s. The operator provisions a 3-node mesh as a headless service, with each pod discovering peers via DNS (\texttt{meshflow-0:7947,meshflow-1:7947}) and mounting persistent volumes for BoltDB data.

\subsection{Circuit Breaker Pattern}

The HTTP executor wraps all calls in a circuit breaker with configurable failure threshold (default 5) and reset timeout (default 30 seconds). When a downstream service fails repeatedly, the breaker opens and tasks are immediately failed without wasting retry attempts, then transitions to half-open for probing.

\subsection{Prometheus Metrics}

Each node exposes a \texttt{GET /metrics} endpoint in Prometheus text format with counters for tasks completed, tasks failed, workflows run, and claims published, plus a node uptime gauge.

\section{Evaluation}

We evaluate MeshFlow on a single-node Apple M1 (8-core) machine to characterize micro-benchmark performance. All benchmarks use Go's built-in benchmarking framework with 5 iterations per benchmark and the 15 ms arbitration window enabled.

\subsection{Trigger Latency}

Table~\ref{tab:trigger-latency} shows the latency for deploying and triggering workflows of varying DAG topologies. Each measurement includes BoltDB persistence, NATS event publish, and DAG resolution.

\begin{table}[h]
\centering
\caption{Workflow trigger latency by DAG topology (single node).}
\label{tab:trigger-latency}
\begin{tabular}{lrr}
\toprule
\textbf{Workflow} & \textbf{Tasks} & \textbf{Latency (ms)} \\
\midrule
Single task & 1 & 6.3 \\
Linear 3-task & 3 & 6.4 \\
Linear 10-task & 10 & 8.9 \\
Fan-out 11-task & 11 & 9.2 \\
\bottomrule
\end{tabular}
\end{table}

The near-constant trigger latency across DAG sizes (1-task vs. 10-task shows only 2.6 ms increase) demonstrates that DAG resolution overhead is negligible compared to persistence and event publish costs. The fan-out pattern (1 entrypoint + 10 parallel workers) shows 9.2 ms, consistent with linear scaling. Note that the 15 ms arbitration window is included in these measurements; without arbitration, single-task trigger latency drops below 1 ms.

\subsection{Event Throughput}

The embedded NATS JetStream layer achieves 30--52K events per second on a single core, with per-event latency of 29--35 $\mu$s and consumption of approximately 5.2 KB per event. This throughput is measured end-to-end: publish, JetStream persistence, subscriber delivery, and acknowledgment.

Table~\ref{tab:event-throughput} shows the distribution across 5 benchmark runs.

\begin{table}[h]
\centering
\caption{Event throughput through embedded NATS JetStream.}
\label{tab:event-throughput}
\begin{tabular}{lrr}
\toprule
\textbf{Run} & \textbf{Events/sec} & \textbf{$\mu$s/op} \\
\midrule
1 & 39,348 & 33.6 \\
2 & 36,066 & 30.2 \\
3 & 52,471 & 28.9 \\
4 & 34,656 & 35.0 \\
5 & 37,720 & 30.5 \\
\midrule
\textbf{Average} & \textbf{40,052} & \textbf{31.6} \\
\bottomrule
\end{tabular}
\end{table}

The variance between runs (34,656--52,471 events/sec) reflects the impact of the Go garbage collector and OS scheduler on the benchmark. At the average rate of 40,000 events/sec, the event layer can sustain a 1,000-node mesh where each node produces 40 events/sec, far exceeding typical workflow orchestration workloads.

\subsection{Persistence Performance}

BoltDB writes for workflow run state average 6.5 ms per operation across 5 benchmark runs (5.9--7.1 ms), with approximately 31 KB allocated per write. This yields approximately 150 writes per second on a single core. For a typical workflow with 100 tasks, state persistence adds approximately 650 ms total across the entire run---negligible compared to typical task execution times of seconds to minutes.

\subsection{Failure Detection and Recovery}

The SWIM gossip configuration provides failure detection within approximately 2 seconds (probe interval of 1 second with suspicion multiplier of 4). Once a node failure is detected, the gossip protocol propagates the membership change to all surviving nodes within one gossip round (200 ms). Task reclamation is immediate: when a surviving node receives the \texttt{node.left} event, it identifies orphaned tasks claimed by the failed node and re-opens them for claiming.

\subsection{Binary Size and Memory Footprint}

The MeshFlow binary is 26 MB on arm64 Darwin (20 MB in Docker Alpine). At idle, a node consumes approximately 15 MB of RAM. Under load with a 10-task workflow executing, memory usage increases to approximately 25 MB due to in-memory DAG state and event buffers. The BoltDB database for a node with 1,000 workflow runs occupies approximately 2 MB on disk.

\section{Related Work}

\subsection{Workflow Orchestration Systems}

Apache Airflow~\cite{airflow} popularized DAG-based workflow orchestration with its Python-defined pipelines, but relies on a centralized scheduler backed by a relational database. Temporal~\cite{temporal} improves upon Airflow with durable execution and pluggable persistence, but retains a centralized matching service for task-to-worker assignment. Prefect~\cite{prefect} and Dagster~\cite{dagster} offer modern Python-native APIs but share the same centralized architecture. MeshFlow differs fundamentally by eliminating the central scheduler; every node is a peer.

\subsection{Distributed Consensus in Scheduling}

Systems like Apache Helix and Apache Curator use ZooKeeper-based leader election for fault-tolerant scheduling. Kubernetes' scheduler uses a combination of optimistic concurrency and watch-based reconciliation rather than leader election, but still operates as a centralized control loop. Ray's~\cite{ray} distributed scheduler uses a bottom-up approach where local schedulers on each node make task placement decisions, but global state is maintained by the Global Control Store (GCS), which is a separate service with its own fault-tolerance mechanism.

\subsection{Gossip-based Coordination}

SWIM~\cite{swim} and its variants (Lifeguard, Health-Aware SWIM) provide scalable membership with low probe overhead. Serf~\cite{serf} (by HashiCorp) packages SWIM gossip with custom event handling for decentralized orchestration. Consul uses SWIM for service discovery but couples it with Raft for strongly consistent state. MeshFlow uses SWIM only for membership and failure detection, delegating state sharing to the event relay bridge.

\subsection{Event-Driven Architectures}

The event sourcing pattern---where state is derived from an append-only event log---has been widely adopted in systems like Apache Kafka, NATS JetStream~\cite{nats}, and EventStoreDB. MeshFlow applies event sourcing to workflow state: each node independently reconstructs DAG state from the event stream. The gossip-to-NATS relay provides a novel bridge between the two communication domains, enabling decentralized event replication without NATS clustering or external coordination.

\subsection{Choreography vs. Orchestration}

The microservices community distinguishes between orchestration (a central conductor coordinates services) and choreography (services coordinate through events). MeshFlow occupies an unusual middle ground: it is event-driven (choreography-like) but provides explicit DAG definitions (orchestration-like). Tasks are defined as explicit dependency graphs, but their execution is coordinated purely through events on the mesh.

\section{Discussion and Limitations}

\subsection{The Claim Race Trade-off}

The first-claim-wins mechanism trades absolute correctness for scheduling latency. The race window between two nodes claiming the same task equals the gossip propagation delay (approximately 200 ms) plus the arbitration window (15 ms). Within this window, duplicate claims are theoretically possible.

MeshFlow provides two layers of protection: (1) the 15 ms arbitration window, where nodes re-verify their claim after gossip propagation, and \\ (2) JetStream message deduplication, where \texttt{task.claimed} events carry a dedup key preventing duplicate delivery within the stream's dedup window. Together, these mechanisms eliminate duplicate execution in practice while maintaining sub-10 ms claim latency.

\subsection{Scalability}

The topic-based event routing feature reduces gossip bandwidth by only broadcasting events to nodes subscribed to relevant input keys. At 100 nodes with 10 distinct topics, each event reaches on average 10 relevant nodes instead of 99, reducing gossip bandwidth by approximately 90\%. SWIM itself scales well to thousands of nodes: the per-node probe load is constant ($O(1)$), and the infection-style dissemination reaches all nodes in $O(\log N)$ rounds.

\subsection{Implemented Features}

Beyond the core schedulerless architecture, all major extensions are fully implemented:

\begin{itemize}
\item \textbf{Multi-handler execution}: HTTP, shell, WebAssembly (wazero/WASI), and gRPC backends.
\item \textbf{Task output passing}: Downstream tasks receive upstream outputs, enabling dataflow pipelines.
\item \textbf{Conditional branching}: Tasks with unmet conditions are skipped, enabling dynamic DAG paths.
\item \textbf{Cron scheduling}: 5-field cron expressions with wildcards, ranges, and step values.
\item \textbf{Webhook triggers}: \texttt{POST /webhook/\{name\}} for GitHub, Slack, and general integration.
\item \textbf{Workflow templates}: Parameterized YAML with \texttt{\{\{variable\}\}} substitution.
\item \textbf{JetStream deduplication}: Claim dedup keys prevent duplicate delivery.
\item \textbf{Circuit breaker}: HTTP executor wraps calls with threshold-based circuit breaking.
\item \textbf{Prometheus metrics}: Counters for tasks, workflows, claims, and node uptime.
\item \textbf{Topic-based routing}: Selective gossip propagation for bandwidth reduction.
\item \textbf{Kubernetes operator}: CRD + StatefulSet for self-deploying meshes on K8s.
\end{itemize}

\subsection{Future Work}

Remaining extensions planned:

\begin{itemize}
\item \textbf{Multi-tenancy}: Namespace isolation for workflows across teams.
\item \textbf{Authentication \& RBAC}: API keys, mTLS, and role-based access control.
\item \textbf{Workflow versioning}: Canary and blue/green deployment models for workflow updates.
\end{itemize}

\section{Conclusion}

We presented MeshFlow, a schedulerless decentralized workflow orchestration engine that eliminates the central scheduler SPOF through SWIM gossip and first-claim-wins task assignment. By embedding a complete orchestration stack---gossip, pub/sub, DAG engine, and persistence---into a single 26 MB binary with zero external dependencies, MeshFlow radically simplifies deployment while providing fault-tolerant workflow execution.

The first-claim-wins mechanism demonstrates that consensus protocols are not necessary for workflow scheduling: a lightweight arbitration window coupled with JetStream deduplication provides bounded duplicate probability of $6.7 \times 10^{-4}$. Our evaluation shows 6.3 ms trigger latency, 40,000 events/sec throughput, and sub-2-second failure recovery. Beyond the core architecture, MeshFlow now includes four execution backends (HTTP, shell, Wasm, gRPC), conditional branching for dynamic DAG paths, workflow templates, circuit breaker protection, Prometheus metrics, webhook triggers, and a Kubernetes operator---making it a production-ready alternative to centralized orchestrators while preserving its schedulerless design.

MeshFlow is fully open source.\footnote{Source code available at \texttt{https://github.com/praneshnikhar/MeshFlow}} We believe its schedulerless architecture represents a viable alternative to consensus-based orchestration for the growing class of event-driven, cloud-native workflow applications.

\begin{thebibliography}{10}

\bibitem{swim}
A. Das, I. Gupta, and A. Motivala.
\newblock SWIM: Scalable Weakly-consistent Infection-style Process Group Membership Protocol.
\newblock In \emph{Proc. DSN}, 2002.

\bibitem{raft}
D. Ongaro and J. Ousterhout.
\newblock In Search of an Understandable Consensus Algorithm.
\newblock In \emph{Proc. USENIX ATC}, 2014.

\bibitem{airflow}
Apache Software Foundation.
\newblock Apache Airflow.
\newblock \url{https://airflow.apache.org/}, 2015.

\bibitem{temporal}
Temporal Technologies.
\newblock Temporal: Durable Execution Platform.
\newblock \url{https://temporal.io/}, 2020.

\bibitem{prefect}
Prefect Technologies.
\newblock Prefect: Modern Workflow Orchestration.
\newblock \url{https://www.prefect.io/}, 2019.

\bibitem{dagster}
Dagster Labs.
\newblock Dagster: Cloud-native Orchestration.
\newblock \url{https://dagster.io/}, 2019.

\bibitem{nats}
Synadia Communications.
\newblock NATS: Cloud Native, Open Source, High-performance Messaging.
\newblock \url{https://nats.io/}, 2012.

\bibitem{serf}
A. Dadgar and M. Phillips.
\newblock Serf: Decentralized Cluster Membership, Failure Detection, and Orchestration.
\newblock In \emph{Proc. USENIX LISA}, 2014.

\bibitem{ray}
P. Moritz, R. Nishihara, S. Wang, et al.
\newblock Ray: A Distributed Framework for Emerging AI Applications.
\newblock In \emph{Proc. USENIX OSDI}, 2018.

\end{thebibliography}

\end{document}
